\section{Theoretical Framework}

The analysis is guided by organizational resilience, dynamic capabilities, and institutional voids theory. Organizational resilience theory frames fragility as limited capacity to absorb disruption and maintain function. Dynamic capabilities theory emphasizes the ability to sense, seize, and reconfigure resources under uncertainty \citep{teece1997dynamic}. Institutional voids theory explains why these firm-level capabilities may be especially consequential in emerging-market settings where finance, infrastructure, skills, and market-support systems are uneven \citep{khanna1997why}.

The reconstructed outcome creates a theoretical puzzle. The official metadata indicate that the operative rule uses innovation, website status, and employment-size information. These variables are often interpreted as signs of capability, modernization, or organizational upgrading. Yet in the recovered screening index they are associated with a higher probability of being fragility-classified. This paper therefore interprets the empirical pattern as a fragility paradox of modernization rather than as evidence that innovation or digital presence causes fragility.

Several mechanisms could explain this paradox. First, innovation may be reactive: firms may introduce new products, services, or processes because they are under competitive or demand pressure, not because they are robust. Under this interpretation, innovation is a symptom of adaptation under stress. Second, growth-stage modernization may create liquidity strain, coordination costs, and managerial overload. A firm that is expanding, digitizing, or changing production processes may face short-run vulnerability even if its long-run prospects improve. Third, institutional incompleteness may limit the returns to digital and process upgrading when finance, infrastructure, skills, and market systems are uneven. Fourth, a website or process change may signal partial modernization rather than full digital maturity. Finally, part of the paradox may be a screening-rule artifact produced by the construction of the dependent variable itself. Future evidence could distinguish these mechanisms by linking the screening index to longitudinal outcomes, investment costs, credit constraints, digital usage intensity, and post-survey firm survival.

This framework leads to a cautious interpretation. The model identifies predictive signals of a constructed fragility classification. It does not show that innovation, website ownership, or employment levels cause fragility. The main theoretical contribution is to show how a capability-oriented screening rule can generate interpretive ambiguity unless variable labels, index construction, and predictive validation are audited together.
